% ============================================================================
\section{概念：x86 特权级模型}
\label{sec:tss-privilege}
% ============================================================================

x86 保护模式提供 4 级保护环（Ring 0 -- Ring 3），实现硬件级别的权限隔离。
实际上大多数操作系统只用两级：

\begin{itemize}
  \item \textbf{Ring 0}（内核态）：可执行所有指令，访问所有内存、I/O 端口
  \item \textbf{Ring 3}（用户态）：只能访问标记为 User 的页，不能执行特权指令（如 \texttt{hlt}、\texttt{lgdt}、\texttt{mov cr0}）
\end{itemize}

\begin{whybox}
为什么需要特权分离？
\begin{enumerate}
  \item \textbf{故障隔离}：用户程序的 bug（空指针、越界写）不会破坏内核数据结构。
  \item \textbf{安全性}：恶意程序无法直接操作硬件（磁盘、网卡），必须通过系统调用请求内核服务。
  \item \textbf{最小特权原则}：每个组件只拥有完成工作所需的最低权限。
\end{enumerate}
\end{whybox}

\begin{figure}[H]
  \centering
  \begin{tikzpicture}[
    ring/.style={draw, thick, circle, minimum size=#1},
  ]
    % Rings (outermost to innermost)
    \node[ring=8cm, fill=blue!5] (r3) {};
    \node[ring=5.8cm, fill=blue!10] (r2) {};
    \node[ring=4cm, fill=blue!18] (r1) {};
    \node[ring=2.2cm, fill=blue!30] (r0) {};

    % Labels inside rings
    \node[align=center] at (0, 0) {\textbf{Ring 0}\\[-2pt]\small 内核};
    \node at (1.2, 1.2) {\small Ring 1};
    \node at (1.7, 1.9) {\small Ring 2};
    \node at (0, 3.5) {\small\textbf{Ring 3} — 用户态};

    % Right side: annotations with enough vertical spacing
    \node[anchor=west, text width=4.5cm, font=\small] at (4.8, 3.0)
      {应用程序、Shell\\只能访问 User 页};
    \node[anchor=west, text width=4.5cm, font=\footnotesize, gray] at (4.8, 1.5)
      {Ring 1/2 几乎不用\\Linux/Windows 只用 0 和 3};
    \node[anchor=west, text width=4.5cm, font=\small] at (4.8, 0)
      {内核代码\\可执行所有指令};

    % Arrow: privilege direction (placed below the rings to avoid overlap)
    \draw[->, thick, red!70!black] (4.8, -1.0) -- (4.8, -2.5)
      node[right, pos=0.5, font=\small] {特权级升高（数值降低）};
  \end{tikzpicture}
  \caption{x86 四级保护环模型}
  \label{fig:protection-rings}
\end{figure}


% ----------------------------------------------------------------------------
\subsection{CPL、DPL、RPL 三级权限检查}
\label{sec:tss-cpl-dpl-rpl}
% ----------------------------------------------------------------------------

x86 有三个相关的特权级字段，理解它们的关系是掌握保护模式的关键：

\begin{description}
  \item[CPL（Current Privilege Level）]
    当前正在执行的代码的特权级。
    存在 \textbf{CS 段寄存器的低 2 位}。
    CPU 每次取指令时都会检查 CPL。

  \item[DPL（Descriptor Privilege Level）]
    段描述符或门描述符中声明的"目标权限级别"。
    表示"要访问这个段/门，调用方至少需要这个特权级"。

  \item[RPL（Requested Privilege Level）]
    段选择子的低 2 位。表示"调用方声称自己的权限"。
    RPL 的作用是防止内核代表用户态去访问不该访问的资源。
\end{description}

\noindent 访问检查规则：

\[
  \max(\text{CPL}, \text{RPL}) \le \text{DPL}
\]

即 CPL 和 RPL 中更低特权的那个（数值更大），不能超过目标的 DPL。

\begin{thinkbox}
为什么需要 RPL？如果只有 CPL 和 DPL 两级检查，存在什么安全隐患？

\emph{提示：考虑用户态通过系统调用请求内核读取一个段。如果没有 RPL，
内核（CPL=0）代替用户去读数据段时，DPL 检查会通过——即使用户本身无权访问。
RPL 让内核可以"以用户身份"去做权限检查。}
\end{thinkbox}

% ----------------------------------------------------------------------------
\subsection{特权级切换的硬件机制}
\label{sec:tss-hw-switch}
% ----------------------------------------------------------------------------

特权级切换有两个方向，机制完全不同：

\subsubsection{Ring 3 → Ring 0（向上切换——中断/异常/系统调用）}

当 CPU 在 Ring 3 执行时遇到中断、异常或 \texttt{int} 指令，会自动执行以下步骤：

\begin{enumerate}
  \item 从 TSS 读取目标特权级的 SS:ESP（即 \texttt{SS0:ESP0}）
  \item 切换到内核栈：加载 SS0 到 SS，ESP0 到 ESP
  \item 将旧的 SS、ESP、EFLAGS、CS、EIP 依次压入\textbf{内核栈}
  \item 如果异常带错误码（如 \#GP、\#PF），额外压入 error code
  \item 从 IDT 门描述符加载新的 CS:EIP（跳转到中断处理程序）
\end{enumerate}

\begin{cpustate}
Ring 3 → Ring 0 切换后：
\begin{itemize}
  \item CPL 从 3 变为 0
  \item SS:ESP 从用户栈切换到 TSS.SS0:ESP0 指向的内核栈
  \item 旧的用户态上下文已安全保存在内核栈上
  \item EFLAGS.IF 被 interrupt gate 自动清零（禁止嵌套中断）
\end{itemize}
\end{cpustate}

\subsubsection{Ring 0 → Ring 3（向下切换——\texttt{iret}）}

\texttt{iret} 原本用于"从中断返回"。它从栈上弹出保存的 EIP/CS/EFLAGS。
如果检测到目标 CS 的 RPL（低 2 位）大于当前 CPL，说明要"返回到更低特权级"，
CPU 会额外弹出 ESP 和 SS：

\begin{enumerate}
  \item 弹出 EIP → 设为新的指令指针
  \item 弹出 CS → 设为新的代码段（RPL=3 → CPL 变为 3）
  \item 弹出 EFLAGS → 恢复标志寄存器
  \item 弹出 ESP → 设为新的栈指针
  \item 弹出 SS → 设为新的栈段
\end{enumerate}

\begin{warnbox}
x86 没有"直接降级到 Ring 3"的指令。\texttt{iret} 是唯一的"向下切换"方式。
我们的技巧是在 Ring 0 手动构造一个"假的中断返回帧"，让 CPU 以为自己在从 Ring 3 的中断中返回。
\end{warnbox}
